\chapter{Introduction}\label{ch:intro}

\section{The question}

Many of the systems on which modern economies depend are feedback systems populated by adaptive agents. In an electronic market, trading algorithms watch prices, trade on what they see, and move the prices they watch. In a power grid, inverters and flexible loads respond to frequency and in doing so change it. In a stream of traffic, each adaptive cruise controller responds to the car in front and becomes the input to the car behind. In each case the operator who is responsible for safety faces agents whose rules are private, heterogeneous and changing.

A natural response is to ask to see inside. Securities regulators require firms to flag algorithmic orders and identify the algorithm behind each decision; the US Consolidated Audit Trail records who stands behind every order; grid operators ask aggregators for device-level models. Behind these requirements sits an implicit claim: that safety cannot be assured without knowing who the agents are.

This thesis tests that claim. Its central question is:
\begin{quote}
\emph{When can a feedback system with unknown adaptive agents be certified safe without identifying the agents, and what does each kind of knowledge cost?}
\end{quote}

\section{The answer in one paragraph}

For the closed-loop systems studied here, the information that certifies safety is generically far smaller and far cheaper than the information that identifies the agents. Prices reveal the agents only through a merged measure of their memory rates and weights (Chapter~\ref{ch:p1}). Recovering the agents from that measure is a moment problem whose cost explodes as agents become similar: the trading capital needed to separate two memory rates $\delta$ apart grows as $\delta^{-4}$, the number of experiments needed to separate $K$ agents as $(\Delta\eta)^{-(4K-2)}$, and with general memory no method does better than accuracy $1/\log N$. Safety needs much less. The worst-case transient loss of a position depends on the agents only through one number, the summed price response $S$, and $S$ is certified at the parametric rate whatever the number and spacing of the agents, provided their memory is bounded (Chapter~\ref{ch:p2}). The framework of Chapter~\ref{ch:framework} explains why: every question an operator can ask is a functional of the merged measure, and its cost is decided by how that functional behaves as memory becomes long. Smooth questions such as ``is the system safe?'' are cheap; questions that need the fine structure of the measure, such as ``who is trading?'', are expensive; and questions that are singular at infinite memory cannot be answered from finite data at all.

\section{Contributions}

\paragraph{Identification limits (Chapter~\ref{ch:p1}).} In a linear cohort model whose reduced form arises from Kalman filtering and regularized online learning, prices identify only the merged rate--weight measure. Separating a known pair costs executed quantity of order $\delta^{-2}$ and capital of order $\delta^{-4}$, with explicit constants and a finite-horizon certificate; $K$ unknown cohorts cost $(\Delta\eta)^{-(4K-2)}$ experiments, matched by Prony's method; general completely monotone memory gives the logarithmic rate; passive data cost the same polynomial order. Class-tagged disclosure removes the barrier and is worth a factor $(\Delta\eta)^{-(4K-4)}$ in samples. Trader-class positions from Binance and the CFTC pin down memories that prices leave unrestricted.

\paragraph{Safety without identification (Chapter~\ref{ch:p2}).} In a limit order book with adaptive cohorts, cumulative transient risk has the closed form $S=\tilde\kappa/(m-q\tilde\kappa)$, a book overshoots if and only if depletion outweighs impact-weighted reversion, and a pre-trade guard on the summed impulse response is the largest safe admission rule. $S$ is certified with $O(\sigma^2\epsilon^{-2}\log(1/\epsilon))$ experiments; a finite window certifies it only under a memory bound; the certified guard is minimax-optimal in lost capacity up to a logarithmic factor; disclosure has a computable value for safe capacity; and reversion is under-supplied in equilibrium. In a pre-registered empirical program the statistic forecasts adverse excursions in thin books, while mechanism tests fail.

\paragraph{A framework that contains both (Chapter~\ref{ch:framework}).} Questions are functionals of the merged measure and observation contracts are noise models on its moments. For linear questions the thesis proves an upper bound on the cost of certification in terms of two quantities of the question, its tail at long memory and its coefficient energy (Theorem~\ref{thm:dich}(i), \Proved{}), and states the matching lower bound as its main open technical target. The rates of Chapters~\ref{ch:p1} and~\ref{ch:p2} appear as cases of one picture.

\paragraph{Beyond linear and stationary agents (Chapter~\ref{ch:beyondlinear}).} A comparison theorem shows that the summed-gain guard stays safe for agents with sector-bounded nonlinear rules (Theorem~\ref{thm:sector}, \Proved{}); a second shows that when agents keep learning, a certificate has a lifetime equal to its safety margin divided by the speed of learning (Theorem~\ref{thm:lifetime}, \Proved{}). Whether the guard survives agents who adapt to it is stated as a conjecture.

\paragraph{Beyond markets (Chapter~\ref{ch:beyondmarkets}).} The framework maps to string stability of vehicle platoons with unknown commercial cruise controllers and to frequency safety in grids with unknown device populations (\Target{}, \Pending{}).

\paragraph{Evidence program (Chapter~\ref{ch:evidence}).} Two tests are specified that would test the thesis directly: a ground-truth test that the certificate computed from prices agrees with the one computed from disclosed classes while the classes themselves are not recoverable from prices, and a prospective test on liquidation data recorded after registration (\Pending{}).

\section{Status of the results}

\begin{table}[h]\centering\small
\begin{threeparttable}
\caption{Main results of the thesis and their status}\label{tab:intro_status}
\begin{tabular}{p{0.52\textwidth}p{0.2\textwidth}p{0.18\textwidth}}
\toprule
Result & Where & Status\\
\midrule
Prices identify only the merged measure & Thm~\ref{p1:thm:merged} & proved\\
Capital of order $\delta^{-4}$; $(\Delta\eta)^{-(4K-2)}$ experiments; $1/\log N$ & Ch.~\ref{ch:p1}, \S\ref{p1:sec:resources}--\ref{p1:sec:mixtures} & proved\\
Safety depends only on $S$; guard is maximal & Lemma~\ref{p2:thm:guard} & proved\\
$S$ certified at the parametric rate; impossible without a memory bound & Thms~\ref{p2:thm:window}, \ref{p2:thm:certcost} & proved\\
Upper bound on certification cost for any linear question & Thm~\ref{thm:dich}(i) & proved here\\
Matching lower bound (the dichotomy) & Thm~\ref{thm:dich}(ii) & target\\
Guard safe for sector-bounded nonlinear agents & Thm~\ref{thm:sector} & proved here\\
Certificate lifetime under learning & Thm~\ref{thm:lifetime} & proved here\\
Guard under strategic adaptation & Conj.~\ref{conj:strategic} & conjecture\\
Platoons and grids & Ch.~\ref{ch:beyondmarkets} & target, pending data\\
Ground-truth and prospective tests & Ch.~\ref{ch:evidence} & pending data\\
\bottomrule
\end{tabular}
\end{threeparttable}
\end{table}

\section{Relation to existing work}

The thesis connects four literatures that rarely meet. \emph{System identification} \citep{t:ljung1999} asks how well a system's parameters can be recovered from input--output data; Chapter~\ref{ch:p1} shows that in closed-loop markets the relevant parameters are zeros of the transfer function and are recoverable only at polynomial or logarithmic rates. \emph{Mixture and moment problems} in statistics give the closest mathematical relatives: the optimal rate for estimating a $k$-atom mixing distribution is $n^{-1/(4k-2)}$ \citep{t:heinrich2018,t:wu2020}, super-resolution has the same Vandermonde geometry \citep{t:candes2014,t:moitra2015}, and the estimation of smooth \emph{functionals} of a distribution can be far easier than estimating the distribution itself \citep{t:lepski1999,t:cai2011}. The dichotomy of Chapter~\ref{ch:framework} is the counterpart of that last observation for safety questions in feedback systems. \emph{Robust and positive control} \citep{t:megretski1997,t:farina2000,t:rantzer2015} supplies the comparison arguments of Chapter~\ref{ch:beyondlinear} and the $\ell_1$ gain that bounds worst-case loss. \emph{Partial identification} in econometrics \citep{t:manski2003} provides the stance: when the structure is not identified, ask which decisions are nonetheless determined. In finance, the closest work is on transient impact and resilience, liquidity spirals and automated-market crashes, reviewed in Chapters~\ref{ch:p1} and~\ref{ch:p2} \citep{t:brunnermeier2009,t:kirilenko2017}.

\section{Reading guide}

Chapter~\ref{ch:framework} can be read first or last: it states the common framework, and each of its results points to the chapter where the special case is proved. Chapters~\ref{ch:p1} and~\ref{ch:p2} are self-contained papers and keep their own notation, appendices and references; Appendix~\ref{app:status} maps the notation across chapters. Chapters~\ref{ch:beyondlinear}--\ref{ch:evidence} build on Chapter~\ref{ch:p2}.
